\chapter{Occam Relational Zero Triad}
\label{chap:relational-zero-triad}

\section{Purpose and proof boundary}

This chapter isolates the smallest declared axiom set needed to formalize the
project's relational notion of zero. It deliberately distinguishes the
non-negative relational action-defect introduced below from the ordinary
mechanical action. In particular, the standard variational equation
\(\delta S=0\) is not identified with the numerical statement \(S=0\).

The result proved here is internal to the declared relational model. It gives
three independent routes to the same half-axis zero. A separate analytic
binding is still required before a zero of \(\xi\) may be identified with a
global relational zero-mode.

\section{Occam axiom set}

\begin{axiom}[Normalized relational complementarity]
\label{ax:rz-complement}
Every admitted local relation \(e\) carries a normalized complementary
coordinate
\[
p_e=(\sigma_e,1-\sigma_e),\qquad 0<\sigma_e<1,
\]
with involution
\[
J_e(\sigma_e)=1-\sigma_e.
\]
Its centered coordinate is
\[
x_e=\sigma_e-\frac12.
\]
\end{axiom}

\begin{axiom}[Positive relational action-defect]
\label{ax:rz-defect}
Let \(E\) be finite or countable. Each relation carries a dynamical vector
\(F_e\) in a positive-definite metric \(G_e\), and positive coefficients
\(w_e,\lambda_e>0\). Define
\[
D_H(\sigma)
=
\ln2-H_2(\sigma)
=
D_{\mathrm{KL}}\!\left(
(\sigma,1-\sigma)
\middle\|
\left(\frac12,\frac12\right)
\right)
\]
and
\[
\boxed{
\mathfrak S_{\rm rel}
=
\sum_{e\in E}
w_e\left[
\|F_e\|_{G_e}^{2}
+
\lambda_e D_H(\sigma_e)
\right]
\in[0,\infty].
}
\]
A local \emph{Lagrange node} is a relation satisfying \(F_e=0\). A global
\emph{relational zero} is a state satisfying
\(\mathfrak S_{\rm rel}=0\).
\end{axiom}

\begin{warning}
The symbol \(\mathfrak S_{\rm rel}\) denotes a non-negative project
action-defect. It is not the usual Lorentzian action
\(S=\int L\,dt\). Consequently the theorem below does not replace the
standard stationarity condition \(\delta S=0\) by \(S=0\).
\end{warning}

\section{First route: complement fixed point}

\begin{theorem}[SOH-RZ001: relational fixed-point zero]
\label{thm:rz-fixed-point}
For every relation admitted by \Cref{ax:rz-complement},
\[
J_e(\sigma_e)=\sigma_e
\iff
\sigma_e=\frac12
\iff
x_e=0.
\]
\end{theorem}

\begin{proof}
The fixed-point equation is
\[
1-\sigma_e=\sigma_e,
\]
hence \(2\sigma_e=1\) and \(\sigma_e=1/2\). By definition
\(x_e=\sigma_e-1/2\), so this is equivalent to \(x_e=0\).
Uniqueness is immediate.
\end{proof}

\section{Second route: Shannon and relative entropy}

\begin{theorem}[SOH-RZ002: information zero]
\label{thm:rz-shannon}
For \(0<\sigma<1\),
\[
D_H(\sigma)\ge0
\]
and
\[
D_H(\sigma)=0
\iff
\sigma=\frac12
\iff
H_2(\sigma)=\ln2.
\]
Moreover \(\sigma=1/2\) is the unique stationary point and strict global
minimum of \(D_H\).
\end{theorem}

\begin{proof}
The identity
\[
D_H(\sigma)
=
D_{\mathrm{KL}}\!\left(
(\sigma,1-\sigma)
\middle\|
(1/2,1/2)
\right)
\]
gives non-negativity and the equality condition
\(D_H=0\iff(\sigma,1-\sigma)=(1/2,1/2)\).
Alternatively,
\[
D_H'(\sigma)=\ln\frac{\sigma}{1-\sigma},
\qquad
D_H''(\sigma)=\frac1\sigma+\frac1{1-\sigma}>0.
\]
Hence the only stationary point is \(\sigma=1/2\), and strict convexity
makes it the unique global minimum.
\end{proof}

\section{Third route: relational Lagrange/action zero}

\begin{theorem}[SOH-RZ003: Relational Lagrange--Zero Theorem]
\label{thm:rz-lagrange}
Under \Cref{ax:rz-complement,ax:rz-defect},
\[
\boxed{
\mathfrak S_{\rm rel}=0
\iff
\forall e\in E:
\quad
F_e=0
\quad\text{and}\quad
\sigma_e=\frac12.
}
\]
Therefore the global relational zero is exactly the simultaneous closure of
all local Lagrange nodes and all complementary half coordinates.
\end{theorem}

\begin{proof}
Every term
\[
w_e\left[
\|F_e\|_{G_e}^{2}+\lambda_eD_H(\sigma_e)
\right]
\]
is non-negative. A finite or countable sum of non-negative terms is zero if
and only if every term is zero. Positive definiteness of \(G_e\) gives
\(\|F_e\|_{G_e}^{2}=0\iff F_e=0\), while
\Cref{thm:rz-shannon} gives
\(D_H(\sigma_e)=0\iff\sigma_e=1/2\). Combining these equivalences proves
the claim.
\end{proof}

For the information potential
\[
V_e(\sigma_e)=\lambda_eD_H(\sigma_e),
\]
the equilibrium equation is independently
\[
\frac{\partial V_e}{\partial\sigma_e}
=
\lambda_e\ln\frac{\sigma_e}{1-\sigma_e}=0,
\]
again selecting only \(\sigma_e=1/2\).

\section{Three-route equivalence}

\begin{corollary}[Relational Zero Triad]
\label{cor:rz-triad}
For every admitted local relation,
\[
\boxed{
x=0
\iff
\sigma=\frac12
\iff
H_2(\sigma)=\ln2
\iff
D_H(\sigma)=0.
}
\]
For the full relational system,
\[
\boxed{
\mathfrak S_{\rm rel}=0
\iff
\bigcap_{e\in E}
\{F_e=0,\ \sigma_e=1/2\}.
}
\]
\end{corollary}

The number of relations and the graph topology do not move the local normalized
complement fixed point. Adding relations adds non-negative constraints; it
does not alter \(1/2\) as the local center.

\section{Spinorial consistency check}

The already established binary-spinor relation
\[
e^{2\pi i\sigma}=-1
\]
has the unique solution \(\sigma=1/2\) in \(0<\sigma<1\). Thus the
spinorial sign supplies an independent consistency check but is not required as
a third axiom.

\section{Boundary with celestial Lagrange points}

The phrase \emph{Lagrange node} in this chapter means a zero of a declared
relational vector field \(F_e\). Standard celestial-mechanics Lagrange
points are equilibria of an effective rotating-frame potential and are not in
general geometric midpoints. No claim that all astronomical Lagrange points
occur at coordinate \(1/2\) is made here.

\section{Riemann binding}

For
\[
s=\sigma+it,
\qquad
\sigma=\Re s,
\]
the functional involution is centered by
\[
x=\Re s-\frac12.
\]
Therefore any state independently established to be a global relational
zero-mode satisfies
\[
x=0
\quad\Longrightarrow\quad
\Re s=\frac12.
\]

\begin{corollary}[Conditional Riemann consequence]
\label{cor:rz-rh}
If every non-trivial zero \(\rho\) of \(\xi\) is independently proved to
be a global relational zero-mode governed by
\Cref{ax:rz-complement,ax:rz-defect}, then
\[
\Re\rho=\frac12
\]
for every such zero.
\end{corollary}

\begin{proof}
Apply \Cref{thm:rz-lagrange} to each admitted zero-mode.
\end{proof}

\status{SOH-RZ001--SOH-RZ003 are exact within the two declared relational
axioms. They close the project's internal three-route zero theorem. The
zeta-to-relational-zero identification in \Cref{cor:rz-rh} is a separate
binding and is not promoted here; the external proof status of RH is
unchanged.}
